\section{Relevant Code}
\subsection{Scale-Adjusting Convolutional Neural Network}
\begin{minted}{python3}
class TPL(nn.Module):
    def __init__(self):
        super(TPL, self).__init__()
    
    def forward(self, a, b, c, x):
        Px = c+(1-c)/(1+torch.exp(-a*(x-b)))
        return Px

class SCNN_TPL3(nn.Module):
    name = "SCNN_TPL3"

    def __init__(self, apriori_relevance, leak_rate, dropout_rate, input_channels):
        super(SCNN_TPL3, self).__init__()

        # Multi-scale Assessing Layers
        self.msal8 = nn.Conv1d(in_channels=1, out_channels=8, kernel_size=8, stride=4)
        self.msal_norm_8 = nn.BatchNorm1d(8)
        self.msal25 = nn.Conv1d(in_channels=1, out_channels=8, kernel_size=25, stride=12)
        self.msal_norm_25 = nn.BatchNorm1d(8)
        self.msal50 = nn.Conv1d(in_channels=1, out_channels=8, kernel_size=50, stride=25)
        self.msal_norm_50 = nn.BatchNorm1d(8)
        self.msal_norm_801 = nn.BatchNorm1d(1)

        # An array of length 801 denoting an apriori understanding of the imporance of each channel
        self.apriori_relevace = apriori_relevance

        # Core CNN Layers
        self.reLU = nn.LeakyReLU(leak_rate)
        self.tanh = nn.Tanh()
        self.dropout = nn.Dropout(p=dropout_rate)

        self.core1 = nn.Conv1d(in_channels=input_channels, out_channels=8, kernel_size=15, stride=7)
        self.core_norm1 = nn.BatchNorm1d(8)

        self.core2_1 = nn.Conv1d(in_channels=8, out_channels=16, kernel_size=9, stride=4)
        self.core_norm2_1 = nn.BatchNorm1d(16)

        self.core3_1 = nn.Conv1d(in_channels=16, out_channels=32, kernel_size=9, stride=4)
        self.core_norm3_1 = nn.BatchNorm1d(32)

        self.core4 = nn.Conv1d(in_channels=32, out_channels=64, kernel_size=2)
        self.core_norm4 = nn.BatchNorm1d(64)

        self.core5 = nn.Linear(1600, 32)
        self.core_norm5 = nn.BatchNorm1d(32)

        self.core6 = nn.Linear(32,16)
        self.core_norm6 = nn.BatchNorm1d(16)

        ### Output Layers
        # Determine a,b,c parameters for TPL
        self.outa = nn.Linear(16,4)
        self.outb = nn.Linear(16,4)
        self.outc = nn.Linear(16,4)
        self.outx = nn.Linear(16,4)

        # Determine output values for each class
        self.tpl = TPL()
    
    def forward(self, x):

        # Multi-Scale assessing Layer
        x_msal8 = self.msal8(x)
        x_msal8 = self.msal_norm_8(x_msal8)
        x_msal25 = self.msal25(x)
        x_msal25 = self.msal_norm_25(x_msal25)
        x_msal50 = self.msal50(x)
        x_msal50 = self.msal_norm_50(x_msal50)
        x_msal_apriori = x * self.apriori_relevace
        x_msal_apriori = self.msal_norm_801(x_msal_apriori)
        
        x_msal_convs = torch.concat([x_msal8, x_msal25, x_msal50], axis=2)
        x_msal_convs_flat = x_msal_convs.view(x_msal_convs.size(0),-1)
        x_msal_apriori_flat = x_msal_apriori.view(x_msal_apriori.size(0),-1)
        x_core = torch.concat([x_msal_convs_flat, x_msal_apriori_flat], axis=1)
        x_core = x_core.view(len(x_core), 1, x_core.size(1))

        # Core CNN Layers
        # First Convolution Layer
        x_core = self.core1(x_core)
        x_core = self.core_norm1(x_core)
        x_core = self.dropout(x_core)
        x_core = self.reLU(x_core)

        # First tanh-relu bock
        x_core = self.core2_1(x_core)
        x_core = self.core_norm2_1(x_core)
        x_core = self.dropout(x_core)
        x_core = self.tanh(x_core)
        x_core = self.reLU(x_core)

        # Second tanh-relu block
        x_core = self.core3_1(x_core)
        x_core = self.core_norm3_1(x_core)
        x_core = self.dropout(x_core)
        x_core = self.tanh(x_core)
        x_core = self.reLU(x_core)

        # Fine convolution layer
        x_core = self.core4(x_core)
        x_core = self.core_norm4(x_core)
        x_core = self.dropout(x_core)
        x_core = x_core.view(len(x_core), -1)

        ## Linear Layers
        x_core = self.core5(x_core)
        x_core = self.core_norm5(x_core)
        x_core = self.dropout(x_core)
        x_core = self.tanh(x_core)
        x_core = self.reLU(x_core)

        x_core = self.core6(x_core)
        x_core = self.core_norm6(x_core)
        x_core = self.dropout(x_core)
        x_core = self.tanh(x_core)
        x_core = self.reLU(x_core)

        # Output Layers
        a = self.outa(x_core)
        b = self.outb(x_core)
        c = self.outc(x_core)
        x_out = self.outx(x_core)
        x_out = self.tpl(a,b,c,x_out)
        
        return x_out
\end{minted}

\subsection{VAE Loss Functions}
\label{sec:code:ELBOs}
\begin{minted}{python3}
# Incorporate the Kullback-Leibler loss to prevent the divergence of probability distributionss
def kl_mse_loss(recon_x, x, mu, logvar, beta=1):
    mse = nn.functional.mse_loss(recon_x, x, reduction='sum') / x.size(0)
    kl = -0.5 * torch.sum(1 + logvar - mu.pow(2) - logvar.exp()) / x.size(0) * beta
    return mse + kl, mse, kl

# The unsupervised ELBO loss function
def elbo_supervised(x, x_rec, mu, logvar, y_logits, y_true, beta=1, label_smoothing=0.25):
    mse = nn.functional.mse_loss(x_rec, x, reduction='sum') / x.size(0)
    kl = -0.5 * torch.sum(1 + logvar - mu.pow(2) - logvar.exp()) / x.size(0)
    ce = nn.functional.cross_entropy(y_logits, y_true, label_smoothing=label_smoothing, 
                                    reduction='sum') / x.size(0)
    return mse + kl*beta, mse, kl, ce

# The supervised ELBO loss function
def elbo_unsupervised(x, model, n_classes=4, beta=1):
    logits, q_y, mu, logvar, hidden = model.infer(x)
    elbo_y = []
    for k in range(n_classes):
        y_k = F.one_hot(torch.full((x.size(0),), k, dtype=torch.long, device=x.device), n_classes).float()
        # recompute mu/logvar conditioned on y_k
        mu_k    = model.fc_mu(torch.cat([hidden, y_k], dim=1))
        logvar_k= model.fc_logvar(torch.cat([hidden, y_k], dim=1))
        z_k     = reparameterise(mu_k, logvar_k, eps_scale=model.eps_scale)
        # reconstruction term
        x_rec_k = model.decoder(torch.cat([z_k, y_k], dim=1))
        mse_k = -F.mse_loss(x_rec_k, x, reduction='none').sum(dim=1)
        # KL term
        kl_k    = -0.5 * torch.sum(1 + logvar_k - mu_k.pow(2) - logvar_k.exp(), dim=1) * beta
        elbo_y.append(mse_k - kl_k)
    elbo_tensor = torch.stack(elbo_y, dim=1)
    # weighted sum and entropy
    expected = torch.sum(q_y * elbo_tensor, dim=1)
    entropy = -torch.sum(q_y * torch.log(q_y + 1e-8), dim=1)
    return torch.mean(expected + entropy)
\end{minted}

\subsection{Reparameterisation Trick}
\label{sec:code:raparam}
\begin{minted}{python3}
# The reparameterisation trick for VAEs
def reparameterise(mu, logvar, eps_scale=1):
    std = torch.exp(0.5 * logvar)
    eps = torch.randn_like(std) * eps_scale
    return mu + eps * std
\end{minted}

% \subsection{Convolutional Encoder - Fully Connected Decoder VAE}
% \begin{minted}{python3}
% class ConvExtractor(nn.Module):
%     def __init__(self, input_dim, hidden_dim, lime_only=False):
%         super().__init__()
%         self.lime_only = lime_only
%         self.relu = nn.ReLU(inplace=True)
%         # normalize raw input
%         self.bn_input = nn.BatchNorm1d(input_dim)
%         # conv tokenizers: output channels=1, so batchnorm on channel dim=1
%         self.conv1 = nn.Conv1d(1, 1, kernel_size=3, stride=2)
%         self.bn1   = nn.BatchNorm1d(1)
%         self.conv7 = nn.Conv1d(1, 1, kernel_size=300, stride=50)
%         self.bn7   = nn.BatchNorm1d(1)
%         # compute sequence lengths
%         seq1 = (input_dim - 3)//2 + 1
%         seq7 = (input_dim - 300)//50 + 1
%         # total flattened length: raw + conv outputs
%         flat_len = input_dim + seq1 + seq7
%         # project to hidden_dim
%         self.fc_hidden = nn.Linear(flat_len, hidden_dim)

%     def forward(self, x: torch.Tensor):
%         # x: (batch, input_dim)
%         x = self.bn_input(x)             # normalize features
%         x = x.unsqueeze(1)               # → (batch,1,input_dim)
%         # conv tokenizers
%         x1 = self.relu(self.bn1(self.conv1(x)))  # → (batch,1,seq1)
%         x7 = self.relu(self.bn7(self.conv7(x)))  # → (batch,1,seq7)
%         # flatten
%         x_raw   = x.squeeze(1)                   # → (batch,input_dim)
%         x1_flat = x1.view(x1.size(0), -1)        # → (batch,seq1)
%         x7_flat = x7.view(x7.size(0), -1)        # → (batch,seq7)
%         feat    = torch.cat([x_raw, x1_flat, x7_flat], dim=1)  # → (batch,flat_len)
%         # hidden representation
%         hidden  = self.relu(self.fc_hidden(feat))
%         return hidden

% # -- Semi-Supervised Convolutional VAE --
% class SemiSupervisedVAECP(nn.Module):
%     name = "SemiSupervisedVAECP"
%     def __init__(self,
%                  input_dim:  int,
%                  hidden_dim: int,
%                  latent_dim: int,
%                  eps_scale:  float = 1.0,
%                  lime_only:  bool = False):
%         super().__init__()
%         self.eps_scale = eps_scale
%         # 1) feature extractor
%         self.extractor = ConvExtractor(input_dim, hidden_dim, lime_only=lime_only)
%         # 2) classifier head: hidden → logits for 4 classes
%         self.classifier_head = nn.Sequential(
%             nn.Linear(hidden_dim, 128),
%             nn.ReLU(inplace=True),
%             nn.Linear(128, 4)
%         )
%         # 3) posterior parameters conditioned on (hidden + one-hot y)
%         self.fc_mu     = nn.Linear(hidden_dim + 4, latent_dim)
%         self.fc_logvar = nn.Linear(hidden_dim + 4, latent_dim)
%         # 4) decoder: reconstruct x from (z + one-hot y)
%         self.decoder = nn.Sequential(
%             nn.BatchNorm1d(latent_dim + 4),
%             nn.Linear(latent_dim + 4, input_dim)
%         )
%         # learned scale & shift + nonlinearity
%         self.scale = nn.Parameter(torch.ones(1))
%         self.shift = nn.Parameter(torch.zeros(1))
%         self.tanh  = nn.Tanh()

%     def infer(self, x, y_true_onehot=None):
%         # extract features & predict class
%         hidden   = self.extractor(x)
%         logits_y = self.classifier_head(hidden)
%         q_y      = torch.softmax(logits_y, dim=-1)
%         # use supervised labels if provided
%         y_enc    = y_true_onehot if y_true_onehot is not None else q_y
%         # condition posterior on class
%         hid_y    = torch.cat([hidden, y_enc], dim=1)
%         mu       = self.fc_mu(hid_y)
%         logvar   = self.fc_logvar(hid_y)
%         return logits_y, q_y, mu, logvar, hidden

%     def forward(self, x, y=None):
%         # y: (batch,) integer labels
%         y_onehot    = None
%         if y is not None:
%             y_onehot = F.one_hot(y, num_classes=4).float()
%         logits_y, q_y, mu, logvar, hidden = self.infer(x, y_onehot)
%         # reparameterisation trick
%         std  = torch.exp(0.5 * logvar)
%         eps  = torch.randn_like(std) * self.eps_scale
%         z    = mu + eps * std
%         # decoder input: z + class encoding
%         y_dec = y_onehot if y_onehot is not None else q_y
%         dec_in = torch.cat([z, y_dec], dim=1)
%         rec    = self.tanh(self.decoder(dec_in)) * self.scale + self.shift
%         return rec, logits_y, mu, logvar

%     def predict(self, x):
%         # return posterior parameters
%         with torch.no_grad():
%             logits_y, q_y, mu, logvar, _ = self.infer(x)
%         return mu, logvar, q_y
% \end{minted}

% \subsection{Convolutional Encoder - Transformer Decoder VAE}
% \begin{minted}{python3}
% class ConvExtractor(nn.Module):
%     def __init__(self, input_dim, hidden_dim, latent_dim, lime_only=False):
%         super().__init__()
%         self.lime_only = lime_only
%         self.relu = nn.ReLU(inplace=True)
%         self.bn_input = nn.BatchNorm1d(input_dim)
%         num_abstract_dims = 1
%         self.input_conv_1 = nn.Conv1d(1, num_abstract_dims, kernel_size=3, stride=2, padding=0)
%         self.bn_conv1 = nn.BatchNorm1d(1201)
%         self.input_conv_7 = nn.Conv1d(1, num_abstract_dims, kernel_size=300, stride=50, padding=0)
%         self.bn_conv7 = nn.BatchNorm1d(43)
%         flattened_length = 3647 #input_dim + num_abstract_dims * ((input_dim - 2) + (input_dim - 299)) #+ (input_dim - 3) + (input_dim - 4) + (input_dim - 5) + (input_dim - 23) + (input_dim - 299))
%         self.fc_hidden = nn.Linear(flattened_length, hidden_dim)

%         # self.bn_fc = nn.BatchNorm1d(hidden_dim)

%     def forward(self, x):
        
%         # print("A")
%         x = self.bn_input(x)
%         x = x.unsqueeze(1)
%         # print("B")
%         x1 = self.input_conv_1(x)
%         x7 = self.input_conv_7(x)
%         x1 = x1.view(x1.size(0), -1)
%         x7 = x7.view(x7.size(0), -1)
%         # print("C")
%         x1 = self.bn_conv1(self.relu(x1))
%         x7 = self.bn_conv7(self.relu(x7))

%         x_convs = torch.concat([x1, x7], dim=1) #x2, x3, x4, x5, x6,
%         x = torch.concat([x.squeeze(1), x_convs], dim=1)
%         hidden = self.fc_hidden(x)
%         hidden = self.relu(hidden)
%         # print("D")
%         # print(x.shape)
%         # print(x)
%         # mu = torch.tanh(self.fc_mu(hidden))+self.fc_mu(hidden)
%         # logvar = torch.tanh(self.fc_logvar(hidden))+self.fc_logvar(hidden)
%         # print("f")
%         # print(f"mu {mu}")
%         # print(f"logvar {logvar}")
%         return hidden
    
% class TransformerDecoder(nn.Module):
%     def __init__(self, latent_dim, d_model, n_heads, n_layers, out_len):
%         super().__init__()
        
%         self.latent_proj = nn.Linear(latent_dim, d_model)

%         # 2) Create learnable query embeddings – one per output position
%         self.query = nn.Parameter(torch.randn(out_len, d_model))

%         # 3) Standard Transformer‑decoder stack (causal self‑attn + cross‑attn)
%         layer = nn.TransformerDecoderLayer(d_model, n_heads, dim_feedforward=4*d_model,
%                                            batch_first=True)
%         self.decoder = nn.TransformerDecoder(layer, num_layers=n_layers)

%         # 4) Final 1×1 convolution (shared across positions)
%         self.readout = nn.Linear(d_model, 1)
%         self.tanh = nn.Tanh()
%         self.scale = nn.Parameter(torch.ones(1)*4.5)

%         mask = torch.triu(torch.ones(out_len, out_len, dtype=torch.bool), 1)
%         self.register_buffer("causal_mask", mask)   # moves with .to(device)

%     def forward(self, z):
        
%         mem   = self.latent_proj(z).unsqueeze(1)      # (B, 1, d_model)
%         query = self.query.unsqueeze(0).expand(z.size(0), -1, -1)  # (B, L, d_model)

%         h = self.decoder(query, mem, tgt_mask=self.causal_mask)
%         return self.tanh(self.readout(h).squeeze(-1))*self.scale

% class SemiSupervisedVAE(nn.Module):
%     name = "CTSemiSupervisedVAE"
%     def __init__(self, input_dim, hidden_dim, latent_dim, d_model, n_heads, n_layers, eps_scale=1, lime_only=False):
%         super().__init__()
%         self.eps_scale = eps_scale
%         self.extractor = ConvExtractor(input_dim, hidden_dim, latent_dim, lime_only=lime_only)
%         self.classifier_head = nn.Sequential(
%             nn.Linear(hidden_dim, 128),
%             nn.ReLU(),
%             nn.Linear(128, 4)
%         )
%         self.fc_mu = nn.Linear(hidden_dim + 4, latent_dim)
%         self.fc_logvar = nn.Linear(hidden_dim + 4, latent_dim)
%         self.decoder = TransformerDecoder(latent_dim + 4, d_model, n_heads, n_layers, input_dim)

%     def infer(self, x, y_true_onehot=None):
%         hidden = self.extractor(x)
%         logits_y = self.classifier_head(hidden)
%         q_y = torch.softmax(logits_y, dim=-1)
%         # Then if supervised, we use the true labels, otherwise use q_y
%         y_encoded = y_true_onehot if y_true_onehot is not None else q_y
%         hidden_and_y = torch.concat([hidden, y_encoded], dim=1)
%         mu = self.fc_mu(hidden_and_y)
%         logvar = self.fc_logvar(hidden_and_y)
%         return logits_y, q_y, mu, logvar, hidden
    
%     def forward(self, x, y=None):
        
%         y_onehot = None
%         if y is not None:
%             y_onehot = torch.nn.functional.one_hot(y, num_classes=4).float()
%         logits_y, q_y, mu, logvar, hidden = self.infer(x, y_onehot)
%         reparam = reparameterise(mu, logvar, eps_scale=self.eps_scale)
%         # If not given labelled data, we use the predicted labels
%         y_decoded = y_onehot if y_onehot is not None else q_y
%         decoder_input = torch.concat([reparam, y_decoded], dim=1)
%         reconstructed_x = self.decoder(decoder_input)
%         return reconstructed_x, logits_y, mu, logvar
    
%     def predict(self, x):
%         reconstructed_x, logits_y, mu, logvar = self.forward(x)
%         return mu, logvar
% \end{minted}

% \subsection{Transformer Encoder - Transformer Decoder VAE}
% \begin{minted}{python3}
% class PositionalEncoding(nn.Module):
%     def __init__(self, d_model, max_len, dropout=0.1):
%         super().__init__()
%         pe = torch.zeros(max_len, d_model)
%         pos = torch.arange(0, max_len, dtype=torch.float).unsqueeze(1)
%         div = torch.exp(torch.arange(0, d_model, 2).float() * -(math.log(10000.0)/d_model))
%         pe[:, 0::2] = torch.sin(pos * div)
%         pe[:, 1::2] = torch.cos(pos * div)
%         self.register_buffer('pe', pe)
%         self.dropout = nn.Dropout(dropout)

%     def forward(self, x):
%         # x: (seq_len, batch, d_model)
%         x = x + self.pe[:x.size(0)].unsqueeze(1)
%         return self.dropout(x)


% class TransformerExtractor(nn.Module):
%     def __init__(self,
%                  input_dim:   int,
%                  hidden_dim:  int,
%                  latent_dim:  int,
%                  lime_only:   bool = False,
%                  num_heads:   int  = 4,
%                  num_layers:  int  = 3,
%                  dropout:     float= 0.1):
%         super().__init__()
%         self.lime_only = lime_only

        
%         self.bn_input = nn.BatchNorm1d(input_dim)

        
%         self.conv1 = nn.Conv1d(1, 1, kernel_size=3,  stride=2)
%         self.bn1   = nn.BatchNorm1d(1)
%         self.conv7 = nn.Conv1d(1, 1, kernel_size=300, stride=50)
%         self.bn7   = nn.BatchNorm1d(1)

%         # compute how many tokens each produces
%         self.seq1 = (input_dim - 3)//2 + 1
%         self.seq7 = (input_dim - 300)//50 + 1
%         self.seq  = self.seq1 + self.seq7

       
%         self.token_proj = nn.Linear(1, latent_dim)

        
%         self.pos_enc = PositionalEncoding(latent_dim, max_len=self.seq, dropout=dropout)
%         layer = nn.TransformerEncoderLayer(
%             d_model=latent_dim,
%             nhead=num_heads,
%             dim_feedforward=hidden_dim,
%             dropout=dropout,
%             activation='relu'
%         )
%         self.transformer = nn.TransformerEncoder(layer, num_layers=num_layers)

%         # 5) pool back to a single vector and project to hidden_dim
%         self.pool    = nn.AdaptiveAvgPool1d(1)
%         self.post_fc = nn.Linear(latent_dim, hidden_dim)
%         self.relu    = nn.ReLU(inplace=True)

%     def forward(self, x: torch.Tensor):
%         # x: (batch, input_dim)
%         x = self.bn_input(x)           
%         x = x.unsqueeze(1)             

%         # conv tokenizers
%         x1 = F.relu(self.bn1(self.conv1(x)))  
%         x7 = F.relu(self.bn7(self.conv7(x)))  

%         # make (seq, batch, 1)
%         tok = torch.cat([x1, x7], dim=2)       
%         tok = tok.permute(2,0,1)              
%         # embed + attend
%         emb = self.token_proj(tok)             
%         emb = self.pos_enc(emb)
%         out = self.transformer(emb)            

%         # pool & project to hidden_dim
%         out = out.permute(1,2,0)               
%         pooled = self.pool(out).squeeze(-1)   
%         hidden = self.relu(self.post_fc(pooled)) 

%         return hidden

% class TransformerDecoder(nn.Module):
%     def __init__(self, latent_dim, d_model, n_heads, n_layers, out_len):
%         super().__init__()
        
%         self.latent_proj = nn.Linear(latent_dim, d_model)

        
%         self.query = nn.Parameter(torch.randn(out_len, d_model))

        
%         layer = nn.TransformerDecoderLayer(d_model, n_heads, dim_feedforward=4*d_model,
%                                            batch_first=True)
%         self.decoder = nn.TransformerDecoder(layer, num_layers=n_layers)

        
%         self.readout = nn.Linear(d_model, 1)
%         self.tanh = nn.Tanh()
%         self.scale = nn.Parameter(torch.ones(1)*4.5)

%         mask = torch.triu(torch.ones(out_len, out_len, dtype=torch.bool), 1)
%         self.register_buffer("causal_mask", mask)   # moves with .to(device)

%     def forward(self, z):
        
%         mem   = self.latent_proj(z).unsqueeze(1)      # (B, 1, d_model)
%         query = self.query.unsqueeze(0).expand(z.size(0), -1, -1)  # (B, L, d_model)

%         h = self.decoder(query, mem, tgt_mask=self.causal_mask)
%         return self.tanh(self.readout(h).squeeze(-1))*self.scale

% class SemiSupervisedVAET(nn.Module):
%     name = "DTSemiSupervisedVAE"
%     def __init__(self, input_dim, hidden_dim, latent_dim, d_model, n_heads, n_layers, eps_scale=1, lime_only=False):
%         super().__init__()
%         self.eps_scale = eps_scale
%         self.extractor = TransformerExtractor(input_dim, hidden_dim, latent_dim, lime_only=lime_only)
%         self.classifier_head = nn.Sequential(
%             nn.Linear(hidden_dim, 128),
%             nn.ReLU(),
%             nn.Linear(128, 4)
%         )
%         self.fc_mu = nn.Linear(hidden_dim + 4, latent_dim)
%         self.fc_logvar = nn.Linear(hidden_dim + 4, latent_dim)
%         self.decoder = TransformerDecoder(latent_dim + 4, d_model, n_heads, n_layers, input_dim)

%     def infer(self, x, y_true_onehot=None):
%         hidden = self.extractor(x)
%         logits_y = self.classifier_head(hidden)
%         q_y = torch.softmax(logits_y, dim=-1)
%         # Then if supervised, we use the true labels, otherwise use q_y
%         y_encoded = y_true_onehot if y_true_onehot is not None else q_y
%         hidden_and_y = torch.concat([hidden, y_encoded], dim=1)
%         mu = self.fc_mu(hidden_and_y)
%         logvar = self.fc_logvar(hidden_and_y)
%         return logits_y, q_y, mu, logvar, hidden
    
%     def forward(self, x, y=None):
        
%         y_onehot = None
%         if y is not None:
%             y_onehot = torch.nn.functional.one_hot(y, num_classes=4).float()
%         logits_y, q_y, mu, logvar, hidden = self.infer(x, y_onehot)
%         reparam = reparameterise(mu, logvar, eps_scale=self.eps_scale)
%         # If not given labelled data, we use the predicted labels
%         y_decoded = y_onehot if y_onehot is not None else q_y
%         decoder_input = torch.concat([reparam, y_decoded], dim=1)
%         reconstructed_x = self.decoder(decoder_input)
%         return reconstructed_x, logits_y, mu, logvar
    
%     def predict(self, x):
%         reconstructed_x, logits_y, mu, logvar = self.forward(x)
%         return mu, logvar
% \end{minted}